% !TEX root = ../main.tex

\chapter{Introdução}
\label{cap:introducao}

Texto da introdução. Cite com \verb|\cite|, \verb|\textcite| ou
\verb|\parencite|: segundo \textcite{doCarmo1992}, \ldots{} \parencite{lima1973variedades}.

Referências cruzadas usam \verb|\cref|/\verb|\Cref|: veja a
\cref{def:exemplo} e o \Cref{teo:exemplo}.

\section{Exemplos de ambientes}

\begin{definition}
  \label{def:exemplo}
  Um \emph{objeto} é \ldots
\end{definition}

\begin{theorem}
  \label{teo:exemplo}
  Todo objeto satisfaz \ldots
\end{theorem}

\begin{proof}
  Demonstração.
\end{proof}

\begin{observation}
  Uma observação.
\end{observation}

\begin{citacao}
  Citação longa (mais de três linhas) no padrão ABNT: recuo de 4\,cm,
  fonte menor e espaçamento simples \parencite[p.~1]{doCarmo1992}.
\end{citacao}
